\begin{claim}{D.28}{Proved}{Prime-order sections and the full fixed subgroup}
Let $p$ be a prime and let $\rho:C_p\to GT_0$ be a continuous
homomorphism whose multiplier has $p$-adic component $1$.
This condition is automatic for odd $p$. Let $\alpha$ be the
normalized automorphism of $F=\widehat F(x,y)$ associated to
the chosen generator of $C_p$, and put
\[
 E=F\rtimes_\alpha C_p,\qquad S=\{1\}\times C_p,
 \qquad C=C_F(S).
\]
There is exactly one $F$-conjugacy class of sections of
$E\to C_p$. Each of the three peripheral normalizers contains
a section. The inclusion $C\hookrightarrow F$ induces an
isomorphism of maximal pro-$p$ quotients
\[
                         C^{(p)}\xrightarrow{\sim}F^{(p)}.
                                                               \tag{D.28a}
\]
This assertion concerns the maximal pro-$p$ quotient of the
actual closed fixed subgroup. It does not identify that
subgroup with $F$. An explicit order-three automorphism below
has a proper fixed subgroup and exactly the isomorphism
(D.28a), while its third peripheral class fails in a finite
dihedral quotient.
\end{claim}
\begin{proof}
The normalized map $GT_0\to\Aut(F)$ is a homomorphism:
for the pair $(\lambda,f)$ it sends
$x\mapsto x^\lambda$, $y\mapsto f^{-1}y^\lambda f$, and
composition gives precisely the product $f\alpha(g)$ in
the second coordinate. Consequently $\alpha^p=1$.
On the maximal pro-$p$ quotient its multiplier is $1$;
the Magnus argument of D.10 and D.25 makes this finite-order
normalized automorphism the identity. Thus the quotient of
the semidirect product is the explicit map
\[
 \pi:E\longrightarrow F^{(p)}\times C_p,
                  \qquad(h,s^j)\longmapsto(\bar h,s^j).
                                                               \tag{D.28b}
\]
For odd $p$, a finite-order unit in $\Z_p^\times$ of order
dividing $p$ is $1$: reduction modulo $p$ eliminates the
prime-to-$p$ part, and the logarithm on $1+p\Z_p$ eliminates
the remaining torsion. This explains the automatic condition.

Write $c_1=x$, $c_2=y$, $c_3=(xy)^{-1}$, and
$L_i=\overline{\langle c_i\rangle}$. Preservation of the
three peripheral classes gives an exact sequence
\[
 1\longrightarrow L_i\longrightarrow N_E(L_i)
                \longrightarrow C_p\longrightarrow1;
                                                               \tag{D.28c}
\]
its kernel is $L_i$ by A.7. Choose $a_i$ over the specified
generator of $C_p$. Then
\[
 a_i c_i a_i^{-1}=c_i^\lambda,\qquad
 a_i^p=c_i^{m_i},\qquad (\lambda-1)m_i=0.
                                                               \tag{D.28d}
\]
The first component of $\pi(a_i)$ centralizes $\bar c_i$,
since $\lambda_p=1$. The centralizer of this primitive
element in the free pro-$p$ group is
$\overline{\langle\bar c_i\rangle}$. Hence
$\pi(a_i)=(\bar c_i^{v_i},s)$ for $v_i\in\Z_p$, and
comparison of $p$th powers yields
\[
                            (m_i)_p=pv_i.
                                                               \tag{D.28e}
\]
Let $N_\lambda=1+\lambda+\cdots+\lambda^{p-1}$.
In the $q$-adic coordinate where $\lambda_q\ne1$,
the equation $\lambda_q^p=1$ and the integral-domain
property give $(N_\lambda)_q=0$ and $(m_i)_q=0$.
Where $\lambda_q=1$, the norm is $p$; division by $p$
is possible for $q\ne p$, and (D.28e) supplies it at
$q=p$. Define $u_i\in\widehat\Z$ coordinatewise by
$(u_i)_q=-(m_i)_q/p$ in the latter coordinates and zero
in the former. This gives
$N_\lambda u_i=-m_i$. The element
$b_i=a_i c_i^{u_i}$ has order $p$, by (D.28d), and
projects to the prescribed generator. Thus every
normalizer in (D.28c) contains a section $S_i$.

Let $e_p\in\widehat\Z$ be the idempotent with
$p$-coordinate $1$ and all other coordinates $0$.
Each $S_i$ centralizes $c_i^{e_p}$. Its image in
$F^{\rm ab}\otimes\F_p$ is the nonzero primitive vector
$[c_i]$. It follows that
\[
       \dim_{\F_p}H^1(C_F(S_i),\F_p)\ge1.
                                                               \tag{D.28f}
\]
This inference uses restriction of a linear functional
on $F^{\rm ab}/p$ that is nonzero on the displayed
element; it does not infer first homology from the mere
presence of a pro-$p$ subgroup.

We now use the finite-cohomology form of the profinite
Brown theorem, quoted as Theorem~2 and applied in
\cite[\S1.1, equations (1)--(6)]{A-KochloukovaZalesskii2007}.
There is no subgroup $C_p\times C_p$ in $E$, because
every finite subgroup injects into $E/F=C_p$.
The Lyndon--Hochschild--Serre spectral sequence and
$\operatorname{cd}_p(F)=1$ give, for every $n>1$,
\[
 \dim_{\F_p}H^n(E,\F_p)\le
             1+\dim_{\F_p}H^1(F,\F_p)=3.
\]
For each subgroup $T$ of order $p$ in $E$, projection
to $C_p$ shows that its normalizer equals its centralizer:
a normalizing element must induce the identity on the
generator of the quotient. Dividing by the unique element
of $T$ with the same quotient gives
\[
                 N_E(T)=C_F(T)\times T.
\]
The Brown restriction map is a dense embedding into the
product of the cohomology of these normalizers. Its
finite-dimensional source makes it an isomorphism and
makes the number of factors finite. Each closed subgroup
$C_F(T)$ has $p$-cohomological dimension at most one.
The direct-product cohomology calculation therefore gives
the exact finite sum and bound
\[
 \sum_{[T]}
    \left(1+\dim_{\F_p}H^1(C_F(T),\F_p)\right)
       =\dim_{\F_p}H^n(E,\F_p)\le3,
                                                               \tag{D.28g}
\]
where $[T]$ runs through $E$-conjugacy classes of order-$p$
subgroups. These are also their $F$-conjugacy classes:
for any such $T$, the decomposition $E=FT$ gives the
same conjugacy orbit under $E$ and $F$.

By (D.28f), two distinct cusp-section classes would
contribute at least four to (D.28g). Thus all three
$S_i$ are conjugate. Conjugating $S_1$ and $S_2$ to a
single representative places conjugates of
$x^{e_p}$ and $y^{e_p}$ in its fixed subgroup.
Their images in $F^{\rm ab}/p$ are still the independent
vectors $[x],[y]$. This gives two independent restricted
linear functionals, so that common factor contributes
at least three. Bound (D.28g) is exhausted. There are
no other classes, and its fixed subgroup has
$\dim H^1(-,\F_p)=2$. Every order-$p$ subgroup projects
isomorphically onto $C_p$ and determines one labelled
section; hence there is exactly one conjugacy class of
sections. In particular these conclusions apply to $S$.

The image of $C$ in $F^{(p)}$ is closed and surjects
on its Frattini quotient by the two independent cusp
vectors, so it is all of $F^{(p)}$. Since $C$ is a
closed subgroup of a free profinite group, it is
projective. Pulling a finite pro-$p$ embedding problem
back to $C$ shows that $C^{(p)}$ is projective pro-$p$,
and hence free pro-$p$. Its rank equals
$\dim H^1(C,\F_p)=2$. The resulting surjection from
a rank-two free pro-$p$ group onto $F^{(p)}$ is an
isomorphism: after identifying free bases it is a
surjective endomorphism of a finitely generated profinite
group, which is injective. This proves (D.28a) in the
claimed direction.

Here is an explicit limitation of that last comparison.
Choose $u\in\widehat\Z^\times$ with $u_q=1$ for
$q\ne7$, and $u_7$ the order-three root with
$u_7\equiv2\pmod7$. The Hensel construction is given
in D.25. The substitutions
\[
                    \beta(x)=x^u,\qquad\beta(y)=y^u
                                                               \tag{D.28h}
\]
have inverse given by exponent $u^{-1}$ and have order
exactly three. The fixed subgroup $C_\beta$ is proper,
since $\beta(x)\ne x$ in the cyclic quotient of order
seven. But it contains $x^{e_3}$ and $y^{e_3}$, so its
map to $F^{(3)}$ is onto. The same Brown bound, now
using this one section and these two independent vectors,
gives $\dim H^1(C_\beta,\F_3)=2$. The preceding
projectivity argument proves
$C_\beta^{(3)}\cong F^{(3)}$, despite $C_\beta\ne F$.

The precise missing cusp condition is detected by
\[
 D_7=\langle r,s\mid r^7=s^2=1, srs=r^{-1}\rangle,
                  \qquad x\mapsto s,\quad y\mapsto sr.
\]
Here $xy\mapsto r$, $u\equiv9\pmod{14}$, and the
two involutions are unchanged under exponent $u$.
Thus $x^u y^u\mapsto r$ whereas $(xy)^u\mapsto r^2$.
The conjugacy class of $r$ consists of $r,r^{-1}$,
so these are not conjugate. Consequently (D.28h) is
not an element of the three-cusp loop-outer group, and
in particular not a $GT_0$ torsion construction. The
program \texttt{code/part\_a/gt\_torsion\_sections.py}
executes this finite certificate and the coordinate
norm-solvability tests used above. D.32 uses the third cusp
in a finite-cover rank argument to strengthen (D.28a) to
$C=F$ under the three-cusp hypotheses, and thereby excludes
the nontrivial prime-order actions considered here.
\end{proof}
